\section{标记}
\subsection{标记线段 \tkzcname{tkzMarkSegment}}
\hypertarget{tms}{}  
  
 \begin{NewMacroBox}{tkzMarkSegment}{\oarg{local options}\parg{pt1,pt2}}% 
该宏用于在线段上放置一个标记。

\medskip
\begin{tabular}{lll}%
\toprule
选项             & 默认值 & 说明   \\
\midrule
\TOline{pos}{.5}{标记的位置} 
\TOline{color}{black}{标记的颜色} 
\TOline{mark}{none}{标记的选择} 
\TOline{size}{4pt}{标记的大小}
\bottomrule
\end{tabular}

可选的标记包括 \TIKZ\ 提供的那些，但也有基于 Yves Combe 的思路所创建的其他标记。
\end{NewMacroBox} 

\subsubsection{多个标记}
\begin{tkzexample}[latex=5cm,small] 
\begin{tikzpicture}
  \tkzDefPoint(2,1){A}
  \tkzDefPoint(6,4){B}
  \tkzDrawSegment(A,B)
  \tkzMarkSegment[color=brown,size=2pt,pos=0.4, mark=z](A,B) 
  \tkzMarkSegment[color=blue,pos=0.2, mark=oo](A,B)
  \tkzMarkSegment[pos=0.8,mark=s,color=red](A,B) 
\end{tikzpicture}
\end{tkzexample}

\subsubsection{\tkzname{mark} 的用法}      
\begin{tkzexample}[latex=5cm,small] 
\begin{tikzpicture}
  \tkzDefPoint(2,1){A} 
  \tkzDefPoint(6,4){B}
  \tkzDrawSegment(A,B)
  \tkzMarkSegment[color=gray,pos=0.2,mark=s|](A,B)
  \tkzMarkSegment[color=gray,pos=0.4,mark=s||](A,B)
  \tkzMarkSegment[color=brown,pos=0.6,mark=||](A,B)
  \tkzMarkSegment[color=red,pos=0.8,mark=|||](A,B)
\end{tikzpicture}
\end{tkzexample}


\subsection{标记多条线段 \tkzcname{tkzMarkSegments}}
\hypertarget{tmss}{} 
 
\begin{NewMacroBox}{tkzMarkSegments}{\oarg{local options}\parg{pt1,pt2 pt3,pt4 ...}}%
参数为由空格分隔的点对列表。可以使用 \TIKZ\ 样式。
\end{NewMacroBox} 

\subsubsection{等腰三角形的标记}      
\begin{tkzexample}[latex=6cm,small]
\begin{tikzpicture}[scale=1]
 \tkzDefPoints{0/0/O,2/2/A,4/0/B,6/2/C}
 \tkzDrawSegments(O,A A,B)
 \tkzDrawPoints(O,A,B)
 \tkzDrawLine(O,B)   
 \tkzMarkSegments[mark=||,size=6pt](O,A A,B)
\end{tikzpicture}
\end{tkzexample} 

\subsection{带构造的标记：等边三角形}   
\begin{tkzexample}[latex=5cm,small] 
 \begin{tikzpicture}[scale=1]
  \tkzDefPoint(0,0){A}\tkzDefPoint(3,2){B} 
  \tkzDefPoint(4,0){C}\tkzDefPoint(2.5,1){P}
  \tkzDrawPolygon(A,B,C)
  \tkzDefEquilateral(A,P) \tkzGetPoint{P'}
  \tkzDefPointsBy[rotation=center A angle 60](P,B){P',C'}
  \tkzDrawPolygon(A,P,P')
  \tkzDrawPolySeg(P',C',A,P,B)
  \tkzDrawSegment(C,P)
  \tkzDrawPoints(A,B,C,C',P,P')
  \tkzMarkSegments[mark=s|,size=6pt,
  color=blue](A,P P,P' P',A) 
  \tkzMarkSegments[mark=||,color=orange](B,P P',C')
  \tkzLabelPoints(A,C) \tkzLabelPoints[below](P) 
  \tkzLabelPoints[above right](P',C',B) 
\end{tikzpicture} 
\end{tkzexample}  

\subsection{标记圆弧 \tkzcname{tkzMarkArc}}
\hypertarget{tma}{}  
  
 \begin{NewMacroBox}{tkzMarkArc}{\oarg{local options}\parg{pt1,pt2,pt3}}% 
该宏用于在圆弧上放置标记：pt1 为圆心，pt2 和 pt3 为圆弧的两个端点。

\medskip
\begin{tabular}{lll}%
\toprule
选项             & 默认值 & 说明   \\
\midrule
\TOline{pos}{.5}{标记的位置} 
\TOline{color}{black}{标记的颜色} 
\TOline{mark}{none}{标记的选择} 
\TOline{size}{4pt}{标记的大小}
\bottomrule
\end{tabular}

可选的标记包括 \TIKZ\ 提供的那些，但也有基于 Yves Combe 的思路所创建的其他标记。
\begin{tkzltxexample}[]
|, ||,|||, z, s, x, o, oo 
\end{tkzltxexample}
\end{NewMacroBox} 

\subsubsection{多个标记}
\begin{tkzexample}[latex=7cm,small] 
\begin{tikzpicture}
\tkzDefPoint(0,0){O}
\pgfmathsetmacro\r{2}
\tkzDefPoint(30:\r){A}
\tkzDefPoint(85:\r){B}
\tkzDrawCircle(O,A)
\tkzMarkArc[color=red,mark=||](O,A,B)
\tkzDrawPoints(B,A,O)
\end{tikzpicture}
\end{tkzexample}

 
\subsection{标记角：\tkzcname{tkzMarkAngle}}
这是一个更为精细的操作，因为选项很多。除 \TIKZ\ 自带的标记符号外，用于标记的符号定义在文件 |tkz-lib-marks.tex| 中，由以下字符指定：\begin{tkzltxexample}[]
|, ||,|||, z, s, x, o, oo 
\end{tkzltxexample}




%                \tkzMarkAngle(B, A, C)
%
% 角标记
% 圆弧（单/双/三）以及相等标记。
%
% 默认值：
%                 arc       = simple
%                 mksize  = 1（圆弧半径）
%                 style traits pleins
%                 mkpos ?  位置：0.5（标记位置）
%                 mark   none
%
% 参数（可选）
%             arc     : l, ll, lll
%             mksize  : 1
%             gap     : 3pt
%             dist    : 1?
%             style   : 线型
%             mkpos   : 0.5
%             mark    : none  , |, ||,|||, z, s, x, o, oo ，不过所有
%  % tikz 符号都是允许的

\begin{NewMacroBox}{tkzMarkAngle}{\oarg{local options}\parg{A,O,B}}%
$O$ 为顶点。可以有多种标记方式：你可以简单地画一段圆弧，或者在该圆弧上添加一个标记。圆弧的样式通过选项 \tkzname{arc} 选择，其半径由 \tkzname{size} 给出，标记的大小由 \tkzname{mksize} 给出；圆弧当然也可以着色。

\medskip

\begin{tabular}{lll}%
\toprule
选项             & 默认值 & 说明                        \\ 
\midrule
\TOline{arc}{l}{选择 l、ll 和 lll（单、双或三重）。}
\TOline{size}{1 (cm)}{圆弧半径。}
\TOline{mark}{none}{标记的选择。}
\TOline{mksize}{4pt}{符号（标记）大小。}
\TOline{mkcolor}{black}{符号（标记）颜色。}
\TOline{mkpos}{0.5}{符号在圆弧上的位置。}
\end{tabular} 
\end{NewMacroBox}  

\DeleteShortVerb{\|}
\subsubsection{\tkzname{mark = x} 与 \tkzname{mark =||} 示例}

\begin{tkzexample}[latex=6cm,small]
\begin{tikzpicture}[scale=.75]
    \tkzDefPoints{0/0/O,5/0/A,3/4/B}
    \tkzMarkAngle[size = 4,mark = x,
                  arc=ll,mkcolor = red,mkpos=.33](A,O,B)
    \tkzMarkAngle[size = 2,mark = ||,
                arc=ll,mkcolor = blue,mkpos=.66](A,O,B)
    \tkzDrawLines(O,A O,B)
    \tkzDrawPoints(O,A,B)
\end{tikzpicture}
\end{tkzexample}

\MakeShortVerb{\|}
\begin{NewMacroBox}{tkzMarkAngles}{\oarg{local options}\parg{A,O,B}\parg{A',O',B'}etc.}%
对于多个角，提供了一个带公共选项的宏。
  \end{NewMacroBox}  

\subsection{标记小角：选项 \tkzname{veclen}}\label{opt-veclen}
  问题来自 "decorate" 操作以及 \tkzcname{tkzMarkAngle} 中 \tkzname{size} 的取值。
  解决方法是将 \tkzcname{tkzMarkAngle} 放置在一个带 \tkzname{veclen} 选项的 scope 中。
  在下例中，若不使用 scope，结果会是：LaTeX Error: Dimension too large。
  
  \begin{tkzexample}[latex=6cm,small]
    \begin{tikzpicture}[scale=1]
      \tkzDefPoint(0,0){O}
      \tkzDefPoint(2.5,0){N}
      \tkzDefPoint(-4.2,0.5){M}
      \tkzDefPointBy[rotation=center O angle 30](N)
      \tkzGetPoint{B}
      \tkzDefPointBy[rotation=center O angle -50](N)
      \tkzGetPoint{A}
      \tkzInterLC[common=B](M,B)(O,B) \tkzGetFirstPoint{C}
      \tkzInterLC[common=A](M,A)(O,A) \tkzGetFirstPoint{A'}
      \tkzDrawSegments(A,C M,A M,B A,B)
      \tkzDrawCircle(O,N)
      \begin{scope}[veclen]
        \tkzMarkAngle[mkpos=.2, size=1.2](C,A,M)
      \end{scope}
      \tkzDrawPoints(O, A, B, M, B, C, A')
      \tkzLabelPoints[right](O,A,B)
      \tkzLabelPoints[above left](M,C)
      \tkzLabelPoint[below left](A'){$A'$}
    \end{tikzpicture}
  \end{tkzexample}

  
\subsection{标记直角：\tkzcname{tkzMarkRightAngle}}

\begin{NewMacroBox}{tkzMarkRightAngle}{\oarg{local options}\parg{A,O,B}}%
选项 \tkzname{german} 会改变标记的样式，选项 \tkzname{size} 会改变其大小。

\medskip
\begin{tabular}{lll}%
\toprule
选项             & 默认值 & 说明         \\ 
\midrule
\TOline{german}{normal}{带内点的德式圆弧。}
\TOline{size}{0.2}{边长。}
\end{tabular} 
\end{NewMacroBox}  

\subsubsection{标记直角} 
\begin{tkzexample}[latex=6cm,small]
\begin{tikzpicture}
  \tkzDefPoints{0/0/A,3/1/B,0.9/-1.2/P}
  \tkzDefPointBy[projection = onto B--A](P)  \tkzGetPoint{H}
  \tkzDrawLines[add=.5 and .5](P,H)
  \tkzMarkRightAngle[fill=blue!20,size=.5,draw](A,H,P) 
  \tkzDrawLines[add=.5 and .5](A,B)
  \tkzMarkRightAngle[fill=red!20,size=.8](B,H,P)
  \tkzDrawPoints[](A,B,P,H)  
\end{tikzpicture}
\end{tkzexample}

\subsubsection{德式直角标记} 
\begin{tkzexample}[latex=6cm,small]
\begin{tikzpicture}
  \tkzDefPoints{0/0/A,3/1/B,0.9/-1.2/P}
  \tkzDefPointBy[projection = onto B--A](P)  \tkzGetPoint{H}
  \tkzDrawLines[add=.5 and .5](P,H)
  \tkzMarkRightAngle[german,size=.5,draw](A,H,P) 
  \tkzDrawPoints[](A,B,P,H) 
  \tkzDrawLines[add=.5 and .5](A,B)
  \tkzMarkRightAngle[german,size=.8](P,H,B) 
\end{tikzpicture}
\end{tkzexample}

\subsubsection{样式混合} 
\begin{tkzexample}[latex=6cm,small]
\begin{tikzpicture}[scale=.75]
  \tkzDefPoint(0,0){A}
  \tkzDefPoint(4,1){B}
  \tkzDefPoint(2,5){C}
  \tkzDefPointBy[projection=onto B--A](C) 
      \tkzGetPoint{H}
  \tkzDrawLine(A,B)
  \tkzDrawLine[add = .5 and .2,color=red](C,H)
  \tkzMarkRightAngle[,size=1,color=red](C,H,A)
  \tkzMarkRightAngle[german,size=.8,color=blue](B,H,C)
  \tkzFillAngle[opacity=.2,fill=blue!20,size=.8](B,H,C)
  \tkzLabelPoints(A,B,C,H)
  \tkzDrawPoints(A,B,C,H)
\end{tikzpicture}
\end{tkzexample}

\subsubsection{完整示例} 

\begin{tkzexample}[latex=6cm,small]
\begin{tikzpicture}[rotate=-90]
\tkzDefPoint(0,1){A}
\tkzDefPoint(2,4){C}
\tkzDefPointWith[orthogonal normed,K=7](C,A)
\tkzGetPoint{B}
\tkzDrawSegment[green!60!black](A,C)
\tkzDrawSegment[green!60!black](C,B)
\tkzDrawSegment[green!60!black](B,A)
\tkzDefSpcTriangle[orthic](A,B,C){N,O,P}
\tkzDrawLine[dashed,color=magenta](C,P)
\tkzLabelPoint[left](A){$A$}
\tkzLabelPoint[right](B){$B$}
\tkzLabelPoint[above](C){$C$}
\tkzLabelPoint[left](P){$P$}
\tkzLabelSegment[auto](B,A){$c$}
\tkzLabelSegment[auto,swap](B,C){$a$}
\tkzLabelSegment[auto,swap](C,A){$b$}
\tkzMarkAngle[size=1,color=cyan,mark=|](C,B,A)
\tkzMarkAngle[size=1,color=cyan,mark=|](A,C,P)
\tkzMarkAngle[size=0.75,color=orange,
    mark=||](P,C,B)
\tkzMarkAngle[size=0.75,color=orange,
   mark=||](B,A,C)
\tkzMarkRightAngle[german](A,C,B)
\tkzMarkRightAngle[german](B,P,C)
\end{tikzpicture} 
\end{tkzexample} 

\subsection{\tkzcname{tkzMarkRightAngles}}
\begin{NewMacroBox}{tkzMarkRightAngles}{\oarg{local options}\parg{A,O,B}\parg{A',O',B'}etc.}%
对于多个角，提供了一个带公共选项的宏。
\end{NewMacroBox}

\subsection{\TIKZ\ 库 \tkzname{angles}} % (fold)
\label{sub:angles_library}

如果你偏好使用 \TIKZ\ 库 \tkzname{angles}，可以通过宏 \tkzcname{tkzPicAngle} 和 \tkzcname{tkzPicRightAngle} 来标记角。

\begin{NewMacroBox}{tkzPicAngle}{\oarg{tikz options}\parg{A,O,B}}%
  
\medskip
\begin{tabular}{lll}%
\toprule
选项             & 示例 & 说明         \\ 
\midrule
\TOline{tikz option}{见下文}{绘制角 $\widehat{AOB}$。}
\end{tabular} 
\end{NewMacroBox}  

\begin{NewMacroBox}{tkzPicRightAngle}{\oarg{tikz options}\parg{A,O,B}}%
  
\medskip
\begin{tabular}{lll}%
\toprule
选项             & 示例 & 说明         \\ 
\midrule
\TOline{tikz option}{见下文}{绘制直角 $\widehat{AOB}$。}
\end{tabular} 

\medskip
\emph{你需要了解库 \tkzname{angles} 的各个选项。}
\end{NewMacroBox} 

\subsubsection{使用 \TIKZ\ 标记角} % (fold)
\label{ssub:angle_with_tikz}


\begin{tkzexample}[latex=7cm,small]
  \begin{tikzpicture}
  \tkzDefPoints{0/0/A,4/0/B}
  \tkzDefTriangle[right,swap](A,B) \tkzGetPoint{C}
  \tkzDrawPolygon(A,B,C)
  \tkzDrawPoints(A,B,C)
  \tkzLabelPoints[below](B,A)
  \tkzLabelPoints[above right](C)
  \tkzPicAngle["$\alpha$",draw=orange,
               <->,angle eccentricity=1.2,
               angle radius=1cm](B,A,C)
  \tkzPicRightAngle[draw,red,thick,
                angle eccentricity=.5,
                pic text=.](C,B,A)
  \end{tikzpicture}
\end{tkzexample}

% subsubsection angle_with_tikz (end)
% subsection angles_library (end)
\endinput
